\documentclass[10pt,a4paper]{amsart}
\usepackage[margin=19mm]{geometry}
\usepackage{amsmath,amssymb,mathtools}
\usepackage[scr=rsfs,cal=euler]{mathalfa}
\usepackage{xfrac}
\usepackage[unicode,hidelinks,bookmarks=false]{hyperref}
\newcommand{\ie}{i.e.\ }
\newcommand{\cf}{cf.\ }
\newcommand{\resp}{respectively\ }
\newcommand{\Subsection}{\subsection}
\providecommand{\nobreakdash}{\nobreak\hbox{-}}
\setlength{\emergencystretch}{2em}
\multlinegap0pt
\usepackage{bm}
\usepackage{bbm}
\usepackage{dsfont}
\usepackage{xspace}
\usepackage{cancel}
\usepackage{mathtools}
\usepackage{amsthm}
\makeatletter
\@ifundefined{theorem}{\newtheorem{theorem}{Theorem}}{}
\makeatother
\newcommand{\N}{\mathbb{N}}	%Atalho
\title[Well-posedness and regularity theory for the fractional diff]{Well-posedness and regularity theory for the fractional diffusion-wave equation in Lebesgue spaces}
\author{Bruno de Andrade, Naldisson Santos}
\date{Source: arXiv:2509.05654v1 (CC BY 4.0)}
\begin{document}
\maketitle

\noindent\emph{Source and licence.} Well-posedness and regularity theory for the fractional diffusion-wave equation in Lebesgue spaces. Version arXiv:2509.05654v1 (CC BY 4.0), distributed under the Creative Commons Attribution 4.0 International licence, as recorded in the arXiv licence metadata for this exact version and shown on the abstract page https://arxiv.org/abs/2509.05654v1. Full rights notice: licenses/arxiv-2509.05654-CC-BY-4.0.txt.\par\smallskip

\noindent\textit{Editorial scope.} Counted unit: the Theorem whose source label is \texttt{Blow} in the article (source0.tex lines 893--894, source label \texttt{Blow}), delivered together with the article's own complete original proof of it (source0.tex lines 895--927, one \texttt{proof} environment of 33 lines). No mathematics is written, repaired or re-derived: statement and proof are the article's own text line for line. Required context retained uncounted: awave (lines 272--277); existence1 (lines 677--685); continuation (lines 859--861). Every definition, display and label of those lines is retained unchanged. Omitted from this package and not redistributed: 1--271, 278--676, 686--858, 862--892, 928--1071. Nothing inside a retained excerpt is omitted or altered: every retained body is delivered exactly as the article's lines are written, so no omission receipt of any kind accompanies this package. Citations: the retained excerpts carry no citation command at all, so no bibliography entry of the article is transcribed and no reference is redistributed. Reference adaptation: none was needed and none was made; every reference inside the delivered unit resolves inside it, so no unresolved reference or citation is produced. Printed slips kept literal: no printed word, symbol, hypothesis or number of the counted statement or of its proof was altered. Layout and class: the article's own class is \texttt{amsart} with packages including amsmath, amssymb, amsfonts, amsthm, geometry, pstricks, pstricks-add, pst-all, pst-math, pst-plot, textcomp, tikz, graphicx, color; the wrapper is the lane's pinned \texttt{amsart} editorial wrapper at 10pt on A4, which restates the theorem-like environments the retained text needs and adds the article's own macro definitions that the retained text needs (\texttt{\textbackslash N}), with extra wrapper packages (bm, bbm, dsfont, xspace, cancel, mathtools). The delivered numbering is the unit's own. Assets: no retained span contains a figure environment, a graphics-inclusion command, a PDF-page inclusion, a table or a TikZ picture; no image file is copied into this package, the package declares no asset dependency and no image file is redistributed with it. Licence: version arXiv:2509.05654v1 of this article is distributed under the Creative Commons Attribution 4.0 International licence, as recorded in the arXiv licence metadata for this exact version and shown on the abstract page https://arxiv.org/abs/2509.05654v1. A full rights notice is delivered at licenses/arxiv-2509.05654-CC-BY-4.0.txt. Characters: every retained excerpt is pure ASCII, so the delivered engine, XeLaTeX with its default Latin Modern text font, reports no missing character.
\begin{eqnarray} \label{awave}
	\left\{\begin{array}{lrr}
		D^{\alpha}_{t}u=A_qu + f(u),\ \ t>0,\\
		u(0)=u_0, \ u'(0)=u_1,\\
	\end{array} \right.
\end{eqnarray}

\begin{theorem}\label{existence1}
	Let $1<\alpha<\frac{2\phi_q}{\pi}$. Consider $\beta \ge 0$ such that $1-\frac{N}{2q'}<\beta<1$ and $1<\rho\leq 1+\frac{2q}{N}(1-\beta )$. Given $v_0\in L^{q}(\Omega)$, we can consider $r>0$ and $\tau >0$ such that for any $u_0, u_1\in B_{r}(v_0)\subset L^{q}(\Omega)$ there exists a unique mild solution $u\in C([0,\tau ], L^{q}(\Omega))$ of the problem $(\ref{awave})$. Furthermore, for all $0\le\theta<\beta$, it follows that
	$$u\in C((0,\tau ],X_{q}^{1+\theta})$$
	and if $\theta>0$ then
	$$\lim_{t\to0}t^{\alpha\theta}\|u(t,u_0,u_1)\|_{X_{q}^{1+\theta}}= 0.$$
	Moreover, if $u_0, w_0, u_1, w_1\in B_r(v_0)\subset L^{q}(\Omega)$, then there exists a constant $c>0$ such that
	$$t^{\alpha\theta}\|u(t,u_0,u_1)-u(t,w_0,w_1)\|_{X_{q}^{1+\theta}}\le c\left(\|u_0 - w_0\|_{L^{q}(\Omega)}+\|u_1 - w_1\|_{L^{q}(\Omega)}\right),$$
	for all $t\in[0, \tau]$.
\end{theorem}

\begin{theorem}\label{continuation}
	Under the conditions of Theorem \ref{existence1}, consider $\tau>0$ and $u: [0,\tau ]\rightarrow X^{1}_{q}$ the mild solution of problem (\ref{awave}). Then, there exist $t_1>\tau$ and a unique continuation $u^*$ of $u$ in $[0,t_1]$.
\end{theorem}

\begin{theorem}\label{Blow} Under the conditions of Theorem \ref{existence1}, let $u$ be the mild solution of problem (\ref{awave}) with maximal time of existence $\tau_{\max}>0$. Then either $\tau_{\max}=\infty$ or $$\displaystyle\limsup_{t\rightarrow \tau_{\max}^-}\|u(t)\|_{X^{1}_{q}}=\infty.$$
\end{theorem}

\begin{proof}
	Suppose that $\tau_{\max}<\infty$ and there exists $k>0$ such that $\|u(t)\|_{X^{1}_{q}}\leq k$, for all $t\in [0, \tau_{\max})$. Let $(t_n)_{n\in\N}\subset[0,\tau_{\max})$ such that $t_n\rightarrow \tau_{\max}$ as $n\rightarrow \infty$. Analogously to the proof of Theorem $\ref{existence1}$, it is shown that 	
	$$\int_{0}^{t_n}(t_n - s)^{\alpha\beta-1}\|\tilde{R}_{\alpha}(t_n-s)f(u(s))-\tilde{R}_{\alpha}(\tau_{\max}-s)f(u(s))\|_{X^{1}_{q}}ds\to0,$$
	as $n\rightarrow \infty$, 
	where $\tilde{R}_{\alpha}(t):= t^{1-\alpha\beta}R_{\alpha}(t)$, $t>0$. 		Suppose $t_n<t_m<\tau_{\max}$. Then,
	\begin{eqnarray*}
		\|u(t_n)-u(t_m)\|_{X^{1}_{q}}&\leq& \|E_{\alpha}(t_n)u_0-E_{\alpha}(t_m)u_0\|_{X^{1}_{q}}\\
		&+& \|S_{\alpha}(t_n)u_1-S_{\alpha}(t_m)u_1\|_{X^{1}_{q}}\\                              
		&+&\int_{0}^{t_n}\|[R_{\alpha}(t_n-s)-R_{\alpha}(t_m-s)]f(u(s))\|_{X^{1}_{q}}ds\\
		&+& \int_{t_n}^{t_m}\|R_{\alpha}(t_m-s)f(u(s))\|_{X^{1}_{q}}ds.
	\end{eqnarray*}
	By strong continuity of $(E_{\alpha}(t))_{t\ge0}$ and $(S_{\alpha}(t))_{t\ge0}$, we have that the two first terms of the right side in the above inequality go to zero as $m, n\rightarrow\infty.$ On the other hand
	\begin{eqnarray*}
		\int_{t_n}^{t_m}\|R_{\alpha}(t_m-s)f(u(s))\|_{X^{1}_{q}}ds &\leq& M \int_{t_n}^{t_m}(t_m-s)^{\alpha\beta-1}\|f(u(s))\|_{X_{\beta}}ds\\
		&\leq& \frac{MC}{\alpha\beta}k^{\rho}(t_m-t_n)^{\alpha\beta}\rightarrow 0,
	\end{eqnarray*}
	as $m,n\rightarrow\infty$. Furthermore, set
	
	$$I=\int_{0}^{t_n}\|[R_{\alpha}(t_n-s)-R_{\alpha}(t_m-s)]f(u(s))\|_{X^{1}_{q}}ds.$$
	Then
	\begin{eqnarray*}
		&&	I \leq \int_{0}^{t_n}[(t_m-s)^{\alpha\beta-1}-(t_n-s)^{\alpha\beta-1}]\|\tilde{R}_{\alpha}(t_n-s)f(u(s))\|_{X^{1}_{q}}ds\\
		&+& \!\!\!\!\!\!\int_{0}^{t_n}(t_n-s)^{\alpha\beta-1}\|\tilde{R}_{\alpha}(t_m-s)f(u(s))-\tilde{R}_{\alpha}(t_n-s)f(u(s))\|_{X^{1}_{q}}ds\\
		&\leq&\frac{MC}{\alpha\beta}k^{\rho}[t_m^{\alpha\beta}-(t_m-t_n)^{\alpha\beta}-t_n^{\alpha\beta}]
	\end{eqnarray*}
	$$+ \int_{0}^{t_n}(t_n - s)^{\alpha\beta-1}\|\tilde{R}_{\alpha}(t_n-s)f(u(s))-\tilde{R}_{\alpha}(\tau_{\max}-s)f(u(s))\|_{X^{1}_{q}}ds$$
	$$+ \int_{0}^{t_m}(t_m - s)^{\alpha\beta-1}\|\tilde{R}_{\alpha}(t_m-s)f(u(s))-\tilde{R}_{\alpha}(\tau_{\max}-s)f(u(s))\|_{X^{1}_{q}}ds$$
	Hence, if $m,n\rightarrow \infty$ it follows that $I\to0$. This computation shows that  $\{u(t_n)\}_{n\in\mathbb{N}}$ is a Cauchy  sequence in $X^{1}_{q}$ and thus there exists $\tilde{u}\in X^{1}_{q}$ such that 
	$\lim_{n\rightarrow\infty}\|u(t_n)-\tilde{u}\|_{X^{1}_{q}}=0.$
	With this, we can extend $u$ to $[0,\tau_{\max}]$ obtaining the equality  
	$$u(t)=E_{\alpha}(t)u_0+S_{\alpha}(t)u_0+ \int_{0}^{t}R_{\alpha}(t-s)f(u(s))ds$$
	for all $t\in[0,\tau_{\max}]$. By Theorem $\ref{continuation}$, we can extend the solution to a bigger interval, which is a contradiction with the maximality of  $\tau_{\max}$.
\end{proof}

\end{document}
